\chapter*{Remerciements}
\addcontentsline{toc}{chapter}{Remerciements}

Je tiens à exprimer mes sincères remerciements à tous ceux qui ont contribué de près ou de loin à la réalisation de ce projet.

En premier lieu, je remercie mon professeur de Sécurité Réseau Avancée pour ses explications claires, ses conseils précieux et ses retours constructifs tout au long du projet. Son expertise dans le domaine de la cybersécurité a été une source d'inspiration et d'apprentissage.

Je remercie également mes camarades de classe pour les discussions techniques enrichissantes et les retours d'expérience échangés pendant la réalisation de ce projet. Ces interactions m'ont permis de mieux comprendre les différentes approches de la sécurisation des systèmes.

Je suis reconnaissant envers la communauté du logiciel libre qui a développé et maintenu les outils utilisés dans ce projet :
\begin{itemize}
    \item L'équipe de développement de \textbf{Linux} et \textbf{Ubuntu}, sans lesquels ce projet n'aurait pas pu être réalisé.
    \item Les contributeurs à \textbf{iptables} et \textbf{netfilter}, qui ont créé un outil puissant et flexible pour la gestion des pare-feu.
    \item L'équipe derrière \textbf{nmap}, qui a développé un excellent outil de scanning réseau.
\end{itemize}

Je remercie également ma famille pour son soutien constant et sa patience pendant les longues heures de travail nécessaires pour compléter ce projet.

Enfin, je remercie toutes les ressources en ligne, les documentation techniques, et les forums de discussion qui m'ont aidé à résoudre les problèmes rencontrés et à approfondir mes connaissances en cybersécurité.

Ce projet a été une expérience enrichissante qui m'a permis d'appliquer concrètement les concepts théoriques appris en classe et de développer une compréhension plus profonde des principes de la sécurité réseau.
